\documentclass[a4paper,10pt]{article}
    \usepackage[utf8]{inputenc} 
    \usepackage[english]{babel}
    \usepackage[scale=.7]{geometry}
    \usepackage{graphicx}
    \usepackage{amsmath,amssymb,amsthm}
    \usepackage{polynomial}
    \usepackage{enumitem}
    \usepackage{tikz}
    \usetikzlibrary{calc}
    \usepackage{hyperref}
    
    \theoremstyle{plain}
    \newtheorem{theorem}{Theorem}[section]
    \newtheorem{lemma}[theorem]{Lemma}
    \newtheorem{proposition}[theorem]{Proposition}
    \newtheorem{corollary}[theorem]{Corollary}
    
    \theoremstyle{definition}
    \newtheorem{definition}[theorem]{Definition}
    \newtheorem{example}[theorem]{Example}
    
    \theoremstyle{remark}
    \newtheorem{remark}[theorem]{Remark}
    
    \title{Definitions}
    
    \author{Johan Voskamp}
    
    \date{\today}
    
    \begin{document}
    \maketitle
    
    \section{}
    
   
    
    \begin{definition}[No-lonely-color graphs]\label{def:lonely}
    	A graph \(G\) is no-lonely-colored if there exists an edge-coloring such that 
        \begin{itemize}
            \item Every vertex's edges can only have at most two of the same color.
            \item Every cycle contains no color exactly once.
        \end{itemize}
    \end{definition}

    \begin{definition}[Cayley Graphs]\label{cayley}
        Let \(G\) be a group, and let \(S
        \subseteq G\), construct the graph \(C(G,S)\) as follows:
        \begin{itemize}
            \item Every \(g\in G\) is a vertex,
            \item For every pair \(s\cdot g_1 = g_2\), where \(s\in S\) and \(g_1, g_2\in G\), set an edge between \(g_1\) and \(g_2\), disallowing identical edges.
        \end{itemize}
        \(C(G,S)\) is then called a Cayley Graph.
    \end{definition}

    \begin{definition}[Minimal Cayley Graphs]\label{min-cayley}
        A Cayley Graph \(C(G,S)\) is called minimal when \(S\) is the smallest possible subset still generating the group \(G\).
    \end{definition}
    
    \begin{definition}[Girth]\label{girth}
        The Girth of a graph is the length of its shortest cycle, if no cycles exist the girth is set to be \(\infty\).
    \end{definition}
    
    \begin{definition}[Multigraphs]\label{multigraph}
        A Multigraph is a graph where it is allowed for two vertices to have multiple edges connecting each other.
    \end{definition}

    \begin{definition}[Bridges]\label{bridge}
        Let \(M\) be a Multigraph and \(e\) be an edge of that graph. If removing \(e\) from \(M\) results in a graph that has more connected components, we call \(e\) a bridge.
    \end{definition}

    \begin{definition}[Hypergraphs]\label{hypergraph}
        Let \(V\) be a set of elements and let \(E\subset P(X)\). We call the collection
        \((V,E),\) an (undirected) Hypergraph. The subset \(E\) is the collection of hyperedges of the Hypergraph.
    \end{definition}

    \begin{definition}[Chromatic Number]\label{def:chromatic}
        The chromatic number of a (hyper)graph \(G\) is the smallest possible number of colors necessary to color the vertices \(G\) if adjacent vertices can't have the same color.
    \end{definition}

    \begin{definition}[r-uniform Hypergraphs]\label{def:r-uniform}
        Let \((V, E)\) be a Hypergraph. If each Hyperedge of the Hypergraph contains exactly \(k\) elements, we call the Hypergraph \(k\)-uniform.
    \end{definition}
    
    \begin{definition}[Tranquil Hypergraphs]\label{def:tranquil}
        ~An \(r\)-uniform hypergraph is tranquil if there exists a set of the same labels that is applied on every hyperedge, it holds that for every possible closed walk in the hypergraph, when projected as a multigraph from label to label, there are no bridges.
    \end{definition}
    
    Geassumeerde voorkennis:
     \begin{itemize}
         \item Grafentheorie
         \item
     \end{itemize}
    \end{document}

         Do I need references for every single definition?
